
\section{1. Introduction}

\subsection{1.1 Motivation}

Autoregressive inference in transformer models is memory-bound, not 
compute-bound. During generation, the model stores keys and values for every 
past token in a \textbf{KV cache}, whose size is:

\[
\text{KV size} = 2 \cdot H_{KV} \cdot d_h \cdot L \cdot B \cdot \text{bytes}
\tag{1.1}
\]

where $H_{KV}$ is the number of KV heads, $d_h$ is the per-head dimension, 
$L$ is the sequence length, and $B$ is the batch size. For a 12-layer GPT-2 
model with $H = 12$, $d_h = 64$, $L = 1024$, and $B = 1$, the KV cache at 
FP16 is:

\[
2 \cdot 12 \cdot 64 \cdot 1024 \cdot 12 \cdot 2 = 37.7 \text{ MB}
\tag{1.2}
\]

which is \textbf{5\% of the model's 774 MB weight footprint} -- small in isolation, 
but growing linearly with sequence length and batch size. At $L = 32{,}768$ 
and $B = 8$, the KV cache dominates inference memory, requiring \textbf{9.6 GB} 
per layer group, exceeding most consumer GPUs.

Two extremes have been explored to reduce this cost:

1. \textbf{Multi-Head Attention (MHA).} Each of the $H$ query heads has its own 
   $K$ and $V$ ($H_{KV} = H$). Maximally expressive, maximally memory-intensive.
2. \textbf{Multi-Query Attention (MQA).} All query heads share a single $K, V$ 
   pair ($H_{KV} = 1$). Minimal memory, but empirically worse quality, 
   especially on tasks requiring fine-grained attention.

\textbf{Grouped-Query Attention (GQA)} was introduced as a middle ground: partition 
the $H$ query heads into $G$ groups, each with a dedicated $K, V$ pair 
($H_{KV} = G$). The KV cache shrinks by a factor of $H / G$, and quality 
degrades gracefully with $G$.

\subsection{1.2 Problem Statement}

Given a pretrained MHA-based model $\mathcal{M}$ with $H$ query heads and 
$H_{KV} = H$, we seek a modified model $\mathcal{M}'$ such that:

\[
\mathcal{M}' \approx \mathcal{M}, \qquad
\text{KV}(\mathcal{M}') \ll \text{KV}(\mathcal{M})
\tag{1.3}
\]

We impose the following constraints:

1. \textbf{Quality preservation:} $\Delta\text{Acc} \leq 5\%$ on SST-2.
2. \textbf{Weight compatibility:} No retraining of base weights; use LoRA recovery.
3. \textbf{Architecture portability:} Must work with flat, nested, and causal transformers.
4. \textbf{Inference ready:} KV cache reduction must be realized at generation time.

\subsection{1.3 Our Approach}

We present a production-hardened implementation of GQA with three key features:

1. \textbf{Exact weight transfer.} The \texttt{copy\_weights} procedure handles four 
   distinct attention configurations: flat (DistilBERT), nested (RoBERTa), 
   causal (GPT-2), and grouped-query (Qwen3.5). For MHA $\to$ GQA conversion, 
   query weights are copied verbatim; KV weights are grouped and repeated.

2. \textbf{Flexible group size.} We support $G \in \{1, 2, \ldots, H\}$. The 
   default is $G = H / 2$, giving a 2$\times$ KV cache reduction.

3. \textbf{LoRA recovery.} After injection, the base weights are frozen and 
   only LoRA adapters are trainable. This preserves pretrained knowledge 
   while allowing task-specific adaptation.

\subsection{1.4 Contributions}

\begin{itemize}
\item \textbf{A production-hardened GQA implementation.} Our PyTorch implementation
\end{itemize}
  handles four attention variants, RoPE, LoRA recovery, and safe fallback.
\begin{itemize}
\item \textbf{A systematic evaluation.} We benchmark GQA on DistilBERT, RoBERTa, and
\end{itemize}
  GPT-2 on SST-2 with ablations over $G / H$ ratio.
\begin{itemize}
\item \textbf{KV cache analysis.} We provide a detailed breakdown of where the memory
\end{itemize}
  savings come from and how they scale with $L$ and $B$.
\begin{itemize}
\item \textbf{Full reproducibility.} Code, checkpoints, and evaluation scripts are
\end{itemize}
  released under MIT.

\subsection{1.5 Paper Organization}

Section 2 reviews related work. Section 3 provides background on attention 
and KV cache. Section 4 presents GQA. Section 5 describes integration with 
four architectures. Section 6 details the experimental setup. Section 7 
reports results. Section 8 analyzes where quality is lost. Section 9 
discusses practical guidance. Section 10 lists limitations. Section 11 
concludes.



\subsection{Notation}
